\chapter{Modelos Locais: Pesos, Memória e uma Execução que Possamos Repetir}
\label{ch:local}

«O modelo ocupa quatro gigabytes; o computador tem oito; deve chegar.» Esta conta ignora tudo o que está à volta do ficheiro. Pode até acertar por sorte. Para planear uma execução, convém saber o que estamos a contar e o que continua fora da conta.

\section{Pesos disponíveis não significam todas as liberdades}
Um modelo com pesos disponíveis permite obter esses parâmetros nas condições da licença. Isso não implica acesso aos dados de treino, código completo, direito de redistribuição irrestrito ou licença open-source. Verifica a licença da revisão concreta, usos permitidos, atribuição e condições de materiais associados.

Um runtime local evita enviar automaticamente pedidos para uma API remota, mas privacidade depende do sistema inteiro: telemetria, extensões, sincronização, logs e ficheiros. Executar localmente também não torna os resultados mais verdadeiros. A vantagem pode ser controlo, disponibilidade ou custo; cada uma precisa de uma comparação própria.

\section{Uma estimativa de memória com hipóteses visíveis}
Para $P$ parâmetros a $b$ bits, a memória ideal de pesos é aproximadamente $Pb/8$ bytes. Com sete mil milhões de parâmetros a quatro bits, obtemos $3.5\times10^9$ bytes: 3.5 GB decimais, cerca de 3.26 GiB. Formatos reais acrescentam escalas, metadados e tensores com outras precisões.

Num descodificador com cache KV densa, uma estimativa ideal é
\[
M_{KV}=2L H_{KV}d_h N B\,\frac{b_{KV}}8.
\]
$L$ é o número de camadas, $H_{KV}$ o de cabeças de keys/values, $d_h$ a dimensão da cabeça, $N$ o número de tokens, $B$ o lote e $b_{KV}$ a precisão da cache em bits. O fator dois conta keys e values. Variantes de arquitetura e de cache podem mudar esta conta.

Num exemplo hipotético com 32 camadas, oito cabeças KV, dimensão 128, 4096 tokens, lote um e cache de 16 bits, o resultado é 536\,870\,912 bytes, ou 512 MiB. O script \texttt{fundamentos.py} executa ambas as contas e os testes conferem os valores. São dimensões didáticas, sem as atribuir a um produto específico.

Ainda faltam ativações, buffers, runtime e memória de outros processos. Para treino faltam também estruturas como estados de otimizador e gradientes. Não concluas que o modelo cabe numa placa de quatro GB porque a soma ideal parece próxima desse número.

\section{A margem desaparece antes do ficheiro}
Vamos manter os valores hipotéticos e assumir uma reserva exata de 4 GiB, isto é, 4294967296 bytes. Os pesos ideais e a cache de 4096 tokens somam 4036870912 bytes. Sobram 258096384 bytes, cerca de 246 MiB, antes de contar as parcelas excluídas. Esta reserva é uma construção matemática, não a memória disponível de um computador ou GPU identificados.

Se duplicarmos o contexto para 8192 tokens nas mesmas condições, a parcela KV duplica para 1073741824 bytes. A soma com os pesos já passa para 4573741824 bytes, acima da reserva hipotética. O ficheiro dos pesos não cresceu. Cresceu o estado necessário à sequência.

O exemplo explica por que «consegui carregar» e «consigo executar a tarefa pretendida» são verificações diferentes. O contexto, a geração e o número de sequências podem mudar a utilização de memória depois do carregamento. Uma aplicação precisa de margem medida sob o perfil de trabalho que vai servir.

O modelo da conta é uma cache densa com as dimensões indicadas. Arquiteturas, formatos, partilha de prefixos ou outras estratégias de runtime podem alterar o comportamento observado. Não devemos transportar a expressão sem verificar essas condições. A função de estimativa torna os pressupostos visíveis; a medição posterior mostra o que a implementação acrescentou ou mudou.

\section{Quantização troca representação, não apenas tamanho}
Quantizar aproxima valores usando uma representação com menos níveis. No exemplo local, uma grelha simétrica de quatro bits usa inteiros de $-7$ a $7$ e uma escala. Os valores $[-1,-0.3,0.2,1]$ são aproximados por $[-1,-2/7,1/7,1]$. O erro de arredondamento nessa grelha fica limitado a meia escala quando não há saturação.

O programa é deliberadamente simples: não empacota bits, não cria GGUF e não reproduz NF4. Mostra a aproximação e permite medir o erro. Num modelo real, a perda de qualidade depende da técnica, tensores, calibração e tarefa. Menos bits não implicam a mesma degradação em todas as capacidades.

Para comparar, conserva base e avaliação e varia uma configuração de cada vez. Mede acerto, memória e tempo. Um formato que reduz memória pode não acelerar um dispositivo sem kernels adequados. O ficheiro pequeno não obrigou o processador a aprender uma instrução nova.

\section{Um erro pequeno no peso pode mudar a decisão}
Considera pesos sintéticos $[1,0.02]$ numa combinação linear e uma entrada $[0,1]$. O score é $0.02$. Na grelha simétrica de quatro bits do nosso exemplo, o maior valor fixa a escala em $1/7$; o peso $0.02$ arredonda para zero. O score da mesma entrada passa a zero.

Se a regra construída decidir pela classe positiva acima de $0.01$, a decisão muda. O erro absoluto num peso foi pequeno e a consequência nesta tarefa foi uma troca de classe. Isto não mede a degradação de um LLM. Mostra por que um erro de representação e um erro na tarefa precisam de avaliações diferentes.

Num sistema real, certas alterações podem não afetar saídas relevantes e outras podem ser amplificadas por operações seguintes ou por decisões perto de um limiar. A comparação com menos bits deve conservar tarefas, critérios e casos difíceis. Um tamanho de ficheiro reduzido é um resultado de armazenamento, não uma certificação das capacidades que interessam.

\section{CPU, GPU e diferentes tempos de espera}
CPUs podem executar modelos locais dentro de limites de memória e tempo. GPUs podem acelerar operações adequadas, com requisitos de drivers, memória e suporte do runtime. Offload parcial distribui trabalho, mas transferências e contenção podem dominar o desempenho.

Separa processamento do contexto, tempo até ao primeiro token e geração seguinte. Uma aplicação que lê um documento longo e responde pouco tem perfil diferente de outra que escreve milhares de tokens. Mede também concorrência: duas sessões podem dividir memória e aumentar a espera mesmo quando cada uma isolada parece rápida.

Os valores publicados por fornecedores ou comunidades precisam de configuração: modelo, quantização, contexto, lote, hardware, versões e comando. Um número de tokens por segundo sem essas condições não permite escolher uma máquina com rigor.

Um ensaio útil mede separadamente a primeira utilização e execuções posteriores quando há carregamento ou caches que alteram o tempo. Também fixa a tarefa: contexto de entrada, comprimento de saída e critérios de paragem. Uma resposta curta pode parecer muito rápida porque fez menos trabalho; um resultado inválido não devia ganhar a comparação por ter terminado primeiro.

Se o objetivo é escolher uma configuração, define primeiro o limite de qualidade e o perfil de recursos aceitável. Só depois compara alternativas dentro desses limites. Este procedimento liga a especialização técnica ao método do primeiro capítulo: a tarefa, a informação e a maneira de confirmar que resultou continuam a decidir a experiência.

\section{Um manifesto para repetir o que aconteceu}
Guarda revisão e hash do modelo, versão do runtime, hardware identificável, contexto, parâmetros, seeds, entradas e saídas. Regista se houve cache, warm-up, paralelismo e ferramentas. Usa pastas novas para experiências novas. Não substituas respostas antigas pelas melhores tentativas.

O ensaio de 03-11-2026 conserva 52 inferências reais com Qwen e llama.cpp~\citep{qwengguf,llamacpp}. A tentativa V3 encontrou uma cópia incompleta do modelo: o runtime recusou o carregamento e o hash diferia da referência. Sem conseguir obter uma cópia íntegra, não houve novo ensaio com LLM. É precisamente o tipo de falha que o manifesto deve tornar visível.

Hashes não demonstram qualidade ou segurança do modelo; demonstram identidade de bytes em relação à referência. Um download que termina sem erro também deve ser conferido quando a integridade faz parte da experiência.

\section{Da experiência local a um serviço}
Abrir um servidor apenas em loopback limita a exposição de rede. Torná-lo acessível a outras máquinas exige autenticação, limites, atualização, monitorização e uma política de dados. Um runtime de laboratório sem chave não se transforma num serviço seguro por mudar o endereço para todas as interfaces.

O adaptador V3 recusa endereços externos e redirecionamentos porque o seu objetivo é um piloto local com dados sintéticos. Não é um cliente universal nem uma plataforma de produção. Os custos de API são zero nesse percurso; eletricidade, equipamento e tempo continuam a existir e não foram convertidos numa fatura imaginária.

\begin{promptlabbox}{O modelo cabe, com que margem?}
\textbf{Objetivo e contexto:} calcular uma estimativa antes de obter pesos. Requisitos: Python 3.10+ ou calculadora; sem GPU e sem rede.

\textbf{Procedimento:} executa o módulo de fundamentos; duplica tokens e depois lote na função de memória; explica quais termos crescem. Lista as parcelas que a função exclui. Se executares depois um modelo, compara o pico observado com esta estimativa ideal.

\textbf{Resultado esperado e validação:} cada duplicação duplica a parcela KV nas hipóteses usadas; o peso ideal mantém-se. O resultado não certifica que um computador concreto suporta a execução.

\textbf{Falhas possíveis:} confundir GB e GiB, contar cabeças de queries em vez de KV, omitir buffers ou apresentar uma previsão de memória como medição de hardware.
\end{promptlabbox}

O tamanho do ficheiro revelou apenas a primeira parcela. Já sabemos procurar cache, buffers, tarefa e identidade dos bytes. A etapa seguinte pergunta como testar e operar o sistema que conseguimos executar, incluindo o dia em que uma atualização faz desaparecer a margem ou altera uma resposta importante.
